\section{결과}
\label{sec:results}

결과는 연구 질문의 순서로 보고하며, 각 절은 해당 그림과 표를 따라 읽도록 구성하였다.
모든 차이는 12개 코호트의 확률효과 통합 추정치와 95\% 신뢰구간으로 나타내며, 괄호 안의 분수는 통합 추정치와 부호가 같은 코호트의 수이다.

\begin{figure}[t]
\centering
\includegraphics[width=\columnwidth]{fig_rq1_effects.pdf}
\caption{단계 구간별 서식지 맥락의 효과}
\label{fig:rq1}
\end{figure}

%======================================================================
\subsection{RQ1: 서식지 맥락과 다음 단계 탈락 위험}
\label{sec:results:rq1}

그림~\ref{fig:rq1}은 네 구성개념이 다음 단계 탈락 위험을 어느 방향으로 바꾸는지를 단계 구간별로 보인다.
모든 단계를 합친 모형에서 주변 경쟁 밀도는 표준편차당 로짓을 $+0.293$ $[0.242, 0.343]$(12/12), 이동 부족분은 $+0.067$ $[0.051, 0.083]$(12/12) 높였고, 남은 팀원 수는 $-0.632$ $[-0.713, -0.551]$(11/11) 낮췄다.
구역 밖 깊이는 방향이 코호트마다 달랐는데(8/12), 이동 부족분과의 상관이 높아 독립된 효과를 가르기 어렵다.

효과의 크기는 단계에 따라 달라졌다.
구역 안쪽 깊이의 보호 효과는 초반에는 뚜렷하지 않았지만(7/12) 중반($-0.289$)과 후반($-0.410$)에는 모든 코호트에서 나타났고, 경쟁 밀도의 효과는 초반 $+0.376$에서 후반 $+0.229$로 작아졌다(구간 간 동일성 검정 $p=5.9\times10^{-5}$).
확률로 바꾸면, 주변의 적 팀이 하나 늘 때 다음 단계 탈락확률은 초반 8.0\%p, 중반 4.1\%p, 후반 3.8\%p 높아졌고, 팀원이 한 명 더 남아 있으면 각각 8.5\%p, 12.6\%p, 14.1\%p 낮아졌다.

이러한 요인은 적은 수의 변수로도 위험을 상당 부분 설명한다.
필요 속도, 주변 경쟁 밀도, 생존 팀 수의 로그, 남은 팀원 수의 네 변수만 쓴 간이 모형은 GBDT가 우연 수준(AUC 0.5)을 넘어 얻은 판별력의 78.8\% $[77.2, 80.3]$를 설명하였다.

%======================================================================
\subsection{RQ2: 기존 설계와의 비교}
\label{sec:results:rq2}

\begin{table*}[t]
\centering
\caption{설계별 예측 성능과 원래 과제 성능}
\label{tab:main}
\footnotesize
\setlength{\tabcolsep}{4pt}
\renewcommand{\arraystretch}{1.15}
\begin{tabular}{lcccccccc}
\hline
 & \multicolumn{6}{c}{본 과제: 다음 단계 탈락} & \multicolumn{2}{c}{KRAFTON의 원래 과제} \\
\cline{2-7}\cline{8-9}
설계 & AUC & 단계 안 AUC & 경기--단계 내부 AUC & 로그 손실 & Brier & 단계별 보정 오차 & 우승 팀 정확도 & C-index \\
\hline
본 연구, GBDT & 0.778 & 0.752 & \textbf{0.757} & \textbf{0.489} & \textbf{0.160} & \textbf{0.013} & 0.351 & 0.654 \\
본 연구, PTR-Former & \textbf{0.779} & \textbf{0.757} & \textbf{0.757} & 0.490 & \textbf{0.160} & 0.017 & 0.386 & 0.663 \\
KRAFTON 적용판 & 0.741 & 0.741 & 0.748 & 0.518 & 0.171 & 0.064 & \textbf{0.417} & \textbf{0.665} \\
KRAFTON + 단계 기저위험 & 0.763 & 0.741 & 0.748 & 0.504 & 0.165 & 0.017 & \textbf{0.417} & \textbf{0.665} \\
Sagong 적용판(XGBoost) & 0.686 & 0.632 & 0.645 & 0.562 & 0.188 & 0.016 & --- & --- \\
\hline
\multicolumn{9}{l}{\emph{같은 KRAFTON 구조에서 학습 목적만 바꾼 경우}} \\
Cox 목적 & 0.742 & 0.731 & 0.745 & 0.517 & 0.171 & 0.056 & \textbf{0.387} & \textbf{0.661} \\
다음 단계 교차엔트로피 목적 & \textbf{0.776} & \textbf{0.755} & \textbf{0.753} & \textbf{0.491} & \textbf{0.161} & \textbf{0.016} & 0.362 & 0.653 \\
\hline
\end{tabular}
\\[2pt]
\parbox{\textwidth}{\scriptsize 12개 코호트 평균. 굵은 글씨는 각 블록에서 가장 좋은 값이다. 우승 팀 정확도와 C-index는 KRAFTON 구현이 정의한 지표이며, 본 연구의 PTR-Former는 순위 출력을 사용하였다.}
\end{table*}

\begin{figure*}[t]
\centering
\includegraphics[width=\textwidth]{fig_contrast.pdf}
\caption{단계별 보정과 설계 요인의 분해}
\label{fig:contrast}
\end{figure*}

표~\ref{tab:main}은 같은 시험 표본에서 설계별 성능을 보인다.
본 연구의 설계는 AUC 0.778(GBDT)과 0.779(PTR-Former)로, KRAFTON 적용판(0.741), 단계별 기저위험을 결합한 KRAFTON(0.763), Sagong 적용판(0.686)보다 높았다.
PTR-Former 기준의 쌍대 차이는 기저위험을 결합한 KRAFTON 대비 $+0.016$ $[0.014, 0.019]$, Sagong 적용판 대비 $+0.089$ $[0.082, 0.096]$였으며, 두 비교 모두 12/12 코호트에서 양수였다.

그림~\ref{fig:contrast}(a)는 이 차이가 어디서 드러나는지 보인다.
실제 다음 단계 탈락률은 단계 1의 19.5\%에서 단계 7의 50.6\%로 높아지지만, KRAFTON 적용판의 평균 예측확률은 26.4\%에서 29.8\%로 거의 평평하다.
이는 상대위험 목적이 단계별 위험 수준을 담지 못한다는 \ref{sec:method:rq2}절의 분석과 일치하며, 단계별 보정 오차는 KRAFTON 0.064, 기저위험 결합 후 0.017, 본 연구의 GBDT 0.013이었다.
차이가 단계 수준에서만 나오는 것은 아니다. 본 연구의 설계는 같은 단계끼리 비교한 단계 안 AUC에서 $+0.017$ $[0.013, 0.021]$(10/12), 같은 경기·같은 순간의 팀끼리 비교한 경기--단계 내부 AUC에서 $+0.010$ $[0.006, 0.015]$(11/12) 높았다.

반대로 KRAFTON의 원래 과제에서는 KRAFTON 적용판이 가장 높았다(표~\ref{tab:main}의 오른쪽 두 열).
우승 팀 정확도는 0.417로 본 연구의 GBDT(0.351)보다 높았고, 같은 KRAFTON 구조 안에서도 Cox 목적은 원래 과제에서, 다음 단계 교차엔트로피 목적은 본 과제에서 더 높은 값을 보였다.
따라서 적용판은 원래 과제에 맞게 학습되었으며, 두 학습 목적은 서로 다른 추정 대상에 맞추어져 있다.

%======================================================================
\subsection{RQ3: 설계 요인의 분해}
\label{sec:results:rq3}

그림~\ref{fig:contrast}(b)는 Sagong 적용판과의 차이를 나눈다.
같은 GBDT에서 Sagong 입력에 서식지 맥락 블록만 더해도 AUC가 0.699에서 $+0.047$ 올랐고, 나머지 입력을 본 연구의 나머지 팀 특징으로 바꾸면 $+0.029$가 더해졌다.
즉 Sagong 설계와의 차이는 대부분 입력 표현, 특히 서식지 맥락 블록에서 온다.

그림~\ref{fig:contrast}(c)는 KRAFTON과의 차이를 나눈다.
입력과 학습 표본을 본 연구에 맞추는 것만으로는 $+0.001$에 그쳤고, Cox 점수에 단계별 기저위험을 결합하면 $+0.016$ $[0.010, 0.022]$, 학습 목적을 다음 단계 교차엔트로피로 바꾸면 $+0.018$ $[0.016, 0.020]$이 더해졌다(모두 12/12).
목적 교체의 이득은 단계 안 AUC($+0.022$)와 경기--단계 내부 AUC($+0.009$)에서도 유지되어, 단계 수준을 맞추는 효과에 그치지 않았다.
반면 명시적 단계 절편($-0.001$), 아키텍처($+0.001$), 단계 축과 관계 편향($+0.002$)은 모두 실질적 동등 범위 안에 있었다.
결과적으로 전체 차이 0.037 가운데 약 0.034가 단계별 기저위험과 학습 목적, 즉 정식화에서 나왔다.



사전 지정한 다섯 비교도 같은 결론을 보였다(부록).
다섯 비교는 모두 실질적 동등으로 판정되었고 REML과 Hartung--Knapp 방식으로 다시 통합해도 판정은 같았으며, 팀 축의 효과는 통계적으로 0보다 컸지만($+0.0027$ $[0.0015, 0.0038]$, Holm $p<0.001$) 그 크기가 0.003에 그쳤다.
학습 예산을 두 배로 늘려도 신경망의 성능은 오르지 않았고(10개 실행 중 9개에서 $-0.0004$에서 $-0.0027$), 순위 보조 손실의 효과도 0.001 이하였다.

마지막으로 입력끼리 비교하면, 본 연구의 입력과 재구성한 KRAFTON 입력은 같은 GBDT에서 예측력이 같았지만($+0.001$ $[-0.002, 0.004]$), 둘을 합치면 $+0.014$ $[0.013, 0.016]$(12/12) 높아져 서로 다른 정보를 담고 있었다. KRAFTON 입력에도 구역까지의 거리와 주변 팀 수가 들어 있으므로(표~\ref{tab:compare}), 두 입력이 담은 정보는 상당 부분 겹칠 수 있다.
가장 큰 두 코호트의 학습 곡선에서는 학습 경기를 모두 써도 AUC가 계속 올랐고, PTR-Former는 전체 학습 경기를 쓸 때에만 GBDT를 넘어섰다(부록).
